% SPDX-License-Identifier: Apache-2.0
% covers: AxiWb
\datasheet{AxiWb}{A bridge from the peripheral end of an AXI4 link to the host lines of a pipelined Wishbone bus.}

\dsfacts{\code{lib/parts/src/bus/wb.rs}}{\code{txhdl\_parts::bus::wb::AxiWb}}%
{\code{//docs:axi}, \code{//docs:vreteno}}

\subsection*{Function}

A controller written elsewhere often speaks Wishbone rather than AXI.
\code{AxiWb} puts such a controller on a link. It is a peripheral
client written as hardware. On one side it holds the channel ends an
\code{AxiPer} gives a client, as one port, \code{bus}, a
\code{PerPort} of the four channels: requests on \code{bus\_req} and write beats on
\code{bus\_w} in, answers on \code{bus\_ans} and read beats on \code{bus\_r} out.
On the other side it drives the host lines of a pipelined Wishbone
bus, \code{cyc}, \code{stb}, \code{we}, \code{adr}, \code{dat} and
\code{sel}, and reads \code{stall}, \code{ack} and the peripheral's
word, \code{rdat}.

It takes one burst at a time and puts each of its words on the lines
as its own request, at consecutive word addresses. The DDR3
controller of Vreteno's board is behind one.

\subsection*{Parameters}

\begin{center}\footnotesize
\begin{tabular}{@{}lp{0.7\textwidth}@{}}
\toprule
Parameter & Meaning \\
\midrule
\code{A} & The link's address width, in bits. \\
\code{I} & The link's identifier width, in bits. \\
\code{AW} & The width of the Wishbone word address, in bits. \\
\bottomrule
\end{tabular}
\end{center}

Words are 32 bits with four lanes, fixed in the source. The tables
below use \code{AxiWb<32, 2, 28>}. There \code{AW} is \code{ddr3::AW},
28 bits, a gigabyte of words, which was the board's DDR3 memory's word
address while this bridge was in front of it.

\dstables{AxiWb}

\subsection*{Behaviour}

\code{stage} says where the request is. The source names five values.

\begin{itemize}
\item \textbf{0, waiting.} A request on \code{bus\_req} is taken when one
offers. The bridge keeps its identifier and its word address, sets
\code{wsel} to all four lanes, and latches \code{left}, the count of
words the burst still has to put on the lines, as \code{len + 1}. A
read goes to stage 2 at once, and a write to stage 1.
\item \textbf{1, a write waiting for its beat.} The beat on \code{bus\_w}
is taken when it offers. Its \code{data} goes to \code{wdat} and its
\code{strb} to \code{wsel}. Then stage 2.
\item \textbf{2, offered.} \code{cyc} and \code{stb} are high. A cycle
with \code{stall} low is the one the peripheral takes the request in,
and the stage moves to 3.
\item \textbf{3, taken.} \code{cyc} alone is high until \code{ack}.
An \code{ack} in the cycle of the take counts too. On \code{ack} the
stage moves to 4, and a read keeps the word on \code{rdat} in
\code{wdat}.
\item \textbf{4, answering.} A read sends a beat on \code{bus\_r}: the
identifier, \code{wdat}, \code{Okay}, and \code{last} high only on
the burst's last word. If words remain the stage returns to 2 at the
next word address with \code{left} one lower; otherwise to 0.
\item A write is answered once, after its last word. Each
acknowledgement before that sends the bridge back to stage 1 for the
next beat, at the next word address. The \code{Answer} of \code{Okay}
on \code{bus\_ans} is sent from stage 4, and each answer waits for room on
its channel.
\item \code{we} is high unless the request is a read. \code{adr},
\code{dat} and \code{sel} are \code{wadr}, \code{wdat} and \code{wsel}.
Every line depends only on the bridge's registers, so nothing
combinational passes from the link to the lines or back.
\item The Wishbone address is a word address: the burst's byte address
shifted down by two and cut to \code{AW} bits. A peripheral decoded at
a region of the link's map sees offsets from the region's base, as
long as the base lies above those bits. A write at
\code{0x4000\_0010} lands in word 4.
\end{itemize}

\subsection*{Verification}

\begin{itemize}
\item Two tests in \code{bus::wb}, in \code{//lib/parts:parts\_test},
run the bridge on a link, behind \code{AxiHost} and \code{AxiPer},
against \code{WbMem}, a simulated Wishbone memory. The first,
\code{a\_read\_gets\_what\_a\_write\_put}, writes and reads one word.
The second, \code{slow\_answers\_and\_a\_calibration}, writes and
reads back 24 words at pseudorandom addresses, against a memory that
stalls for its first 40 cycles and answers after 5.
\item \code{ex\_axi\_wb} writes three words and reads them back, against
a memory that stalls for 10 cycles and answers after 2. The build
simulates the netlist of \code{AxiWb<32, 2, 28>} against that run under
nvc and Verilator: \code{//docs:axi\_wb\_sim\_axi\_wb\_tb\_test} and
\code{//docs:axi\_wb\_vsim\_test}.
\end{itemize}

\subsection*{Used by}

The example \code{ex\_axi\_wb}, and
\code{a\_burst\_is\_served\_word\_by\_word} in the module's own
tests, which writes sixteen words and reads them back as one burst
each way. Sixteen is what \code{LineFetch} issues by default.
Until issue~\#1023 put the board's DDR3 memory behind the controller's
own AXI4 port, \code{Ddr3Per} held this bridge in front of a Wishbone
wrapper of the controller.
That test asserts the COUNT of beats and not only their values: a
bridge that answered one beat would return the right word in it, so
checking the words alone passes on the fault.

\subsection*{Impact of the single-beat bridge (prior to issue 471)}

This behavior is documented here because joining a burst-issuing client
to a peripheral that returns fewer beats than requested risks protocol
failure that neither engine can detect.

\code{LineFetch} decides a burst is finished by its own beat counter
and never reads the returned beat's \code{last}. Against a bridge
that answered one beat it took that word and waited for the rest for
ever, holding \code{running} high: sixty four words became one, and
the run did not end.

\code{LineStore} fared worse and quietly. A write burst is finished
by its response, so the bridge answered after one beat while the rest
of the burst sat in the channel: nine words of sixty four landed and
the rest of memory was left as it was, with nothing to say why.

\subsection*{Limits}

It serves a burst one request at a time rather than as a pipelined
run of several. That removes a link round trip per word and not the
word itself: a burst of sixteen is sixteen Wishbone request and
acknowledge handshakes, one after another. Issuing several requests
before their acknowledgements would remove more, and wants somewhere
to hold answers the link is not ready for, which the bridge does not
have. It always answers \code{Okay}, and it reads no
error line from the peripheral. Words are 32 bits with four lanes.
Address bits above the \code{AW} bits it keeps are dropped.
